\daytitle{25}{光学基础}

\begin{learningbox}{本节导读}
本节先建立“光线模型”：反射、折射、折射率和全反射；再通过色散、双缝干涉、衍射和偏振，把光从“几何光学”推进到“波动光学”。学习顺序是：先会画光路图，再由几何关系写公式，最后理解光在适当条件下可近似沿直线传播，也具有波动性。
\end{learningbox}

\textbf{本节读题路线：} 几何光学题先画界面和法线，再标角度；折射题先判断光从光疏到光密还是光密到光疏，预判靠近或远离法线，再用折射定律计算；色散题比较不同颜色光的折射率；干涉题先确认两束光是否相干，再看条纹间距受 $L$、$\lambda$、$d$ 哪个量影响。

\section*{一、学习目标}
\begin{enumerate}
  \item 能画出反射、折射光路，并使用折射定律和折射率判断光线偏折。
  \item 会先判断光从光疏到光密还是从光密到光疏，再预判光线靠近或远离法线。
  \item 能判断全反射条件，并使用临界角公式。
  \item 能说明色散现象与不同颜色光折射率不同的关系，并能联系光谱分析的基本思想。
  \item 能说明光的频率、波长、波速和颜色之间的基本关系。
  \item 能用双缝干涉公式判断条纹间距与波长、缝距、屏距的关系，并能先统一单位再计算。
  \item 能用典型现象说明干涉、衍射和偏振体现的光的波动性。
\end{enumerate}
\section*{二、初高中衔接}

初中学习光学时，重点是反射、折射和基本成像现象。高中要更精确地处理三个问题：角度必须相对于法线来量；折射率反映光速在介质中的变化；光路图不仅是画图题，也是建立几何关系的工具。

\textbf{科学发展视角：} 光学不只服务于日常照明。光纤通信依赖全反射，光谱分析依赖色散，芯片光刻和精密测量要考虑波长、干涉和衍射。学习本节时，先用光线模型处理折射和全反射，再用波动模型解释干涉、衍射和分辨率。

\section*{三、核心概念}

\begin{conceptbox}{反射与折射}

光遇到两种介质的分界面时，一部分返回原介质形成反射光，一部分进入另一介质形成折射光。

\[
\boxed{i=i'},\qquad
\boxed{n_1\sin i=n_2\sin r}.
\]

\begin{center}
% TikZ 图：反射和折射
% Day 25：光学基础
\begin{tikzpicture}[scale=0.9,>=Latex]
  \draw[thick] (-3.0,0) -- (3.0,0);
  \node[right] at (3.0,0) {界面};
  \draw[dashed] (0,-2.0) -- (0,2.2) node[above] {法线};
  \draw[->,blue,thick] (-2.2,1.7) -- (0,0);
  \draw[->,blue,thick] (0,0) -- (2.0,1.55);
  \draw[->,red,thick] (0,0) -- (1.05,-1.8);
  \draw (126:0.62) arc (126:90:0.62);
  \node at (106:0.9) {$i$};
  \draw (90:0.78) arc (90:38:0.78);
  \node at (62:1.05) {$i'$};
  \draw (-90:0.78) arc (-90:-60:0.78);
  \node at (-75:1.05) {$r$};
  \node[blue] at (-2.0,1.95) {入射光};
  \node[blue] at (2.1,1.8) {反射光};
  \node[red] at (1.35,-1.75) {折射光};
  \node[font=\small] at (-1.8,-0.35) {$n_1$};
  \node[font=\small] at (-1.8,0.35) {空气};
  \node[font=\small] at (2.1,-0.35) {$n_2$};
\end{tikzpicture}


\figcaption{反射和折射：入射角、反射角、折射角都相对于法线测量}\label{fig:reflection-refraction}
\end{center}

\begin{itemize}
  \item $i$：入射角；$i'$：反射角；$r$：折射角。
  \item 光从光疏介质进入光密介质，折射光线靠近法线。
  \item 光从光密介质进入光疏介质，折射光线远离法线。
\end{itemize}

\textbf{术语提醒：} “光疏”和“光密”是对同一频率的光比较两种介质时的相对说法；折射率较小的介质称为光疏介质，折射率较大的介质称为光密介质。

\end{conceptbox}

\begin{conceptbox}{光路可逆性}

在同一组静止介质和光学元件中，如果光可以沿某条路径从 $A$ 到 $B$，那么从 $B$ 沿相反方向入射，也会沿原路径返回 $A$。这叫\textbf{光路可逆性}。

\begin{center}
\begin{tikzpicture}[scale=0.92,>=Stealth]
  \tikzset{
    ray/.style={very thick,blue,->},
    reverse/.style={very thick,orange!90!black,->},
    guide/.style={gray,dashed},
    medium/.style={fill=blue!5,draw=none}
  }

  \fill[medium] (-3.4,-1.2) rectangle (3.4,0);
  \draw[thick] (-3.4,0) -- (3.4,0);
  \draw[guide] (0,-1.20) -- (0,1.45);
  \node[font=\small] at (-2.65,0.42) {空气};
  \node[font=\small] at (-2.65,-0.70) {玻璃};
  \node[font=\scriptsize,gray] at (0.34,1.18) {法线};

  \coordinate (A) at (-2.75,1.05);
  \coordinate (O) at (0,0);
  \coordinate (B) at (2.75,-0.86);

  \draw[very thick,blue] (A) -- (O) -- (B);
  \draw[ray] (-1.72,0.66) -- (-1.22,0.47);
  \draw[ray] (1.10,-0.34) -- (1.60,-0.50);
  \draw[reverse] (1.72,-0.54) -- (1.22,-0.38);
  \draw[reverse] (-1.10,0.42) -- (-1.60,0.61);

  \fill[red] (A) circle (1.5pt) node[above left,font=\small] {$A$};
  \fill[red] (B) circle (1.5pt) node[below right,font=\small] {$B$};
  \fill[black] (O) circle (1.2pt);

  \draw[guide] (0.72,0) arc[start angle=0,end angle=159,radius=0.72];
  \node[font=\scriptsize] at (-0.46,0.30) {$i$};
  \draw[guide] (0.55,0) arc[start angle=0,end angle=-17,radius=0.55];
  \node[font=\scriptsize] at (0.54,-0.16) {$r$};

  \node[font=\small,align=center,blue] at (-1.55,-1.62) {从 $A$ 到 $B$};
  \node[font=\small,align=center,orange!90!black] at (1.55,-1.62) {反向仍沿原路};
  \node[font=\small,align=center] at (0,-2.13) {同一组介质中，光路可逆；作图时可从物点画向像点，也可反向检验};
\end{tikzpicture}


\figcaption{光路可逆性：同一条折射光路反向传播时仍满足相同的折射定律}\label{fig:optical-reversibility}
\end{center}

\textbf{学习意义：} 画反射、折射和光纤传光图时，可以用光路可逆性检查图像是否合理。例如一条折射光线若从 $A$ 到 $B$ 画对了，反向从 $B$ 出发也应沿同一路径回到 $A$。

\textbf{模型边界：} 光路可逆性讨论的是静止、线性、互易介质中的传播路径，而不是光强一定保持不变。普通吸收会使反向传播的光变弱，却不必改变几何路径；介质随时间变化或使用非互易光学器件时，光路才可能不再满足这种简单的可逆关系。高中阶段通常只在反射、折射和光纤等理想模型中使用。

\end{conceptbox}

\begin{conceptbox}{折射率}

介质的绝对折射率定义为
\[
\boxed{n=\frac{c}{v}},
\]
其中 $c$ 是真空中光速，$v$ 是光在介质中的传播速度。

\begin{itemize}
  \item $n$ 越大，光在该介质中传播越慢。
  \item 真空折射率为 1，空气折射率通常近似取 1。
  \item 折射率要针对指定介质和指定频率来讨论；折射角由两种介质的折射率和入射角共同决定。
\end{itemize}

\begin{center}
\begin{tikzpicture}[scale=0.95,>=Stealth]
  \fill[blue!4] (-3.1,0) rectangle (3.1,1.75);
  \fill[orange!8] (-3.1,-1.75) rectangle (3.1,0);
  \draw[thick] (-3.1,0) -- (3.1,0);
  \draw[dashed] (0,-1.6) -- (0,1.6);
  \node[font=\small] at (-2.05,1.42) {空气：$n_1$ 小，$v_1$ 大};
  \node[font=\small] at (-2.05,-1.42) {玻璃：$n_2$ 大，$v_2$ 小};

  \draw[blue,very thick,->] (-2.25,1.25) -- (-0.12,0.07);
  \draw[blue,very thick,->] (0.12,-0.07) -- (0.95,-1.25);

  \draw[gray,thin] (-1.86,0.73) -- (-1.50,1.38);
  \draw[gray,thin] (-1.34,0.44) -- (-0.98,1.09);
  \node[font=\scriptsize,gray] at (-1.63,0.48) {$\lambda_1$};
  \draw[gray,thin] (0.22,-0.42) -- (0.82,0.00);
  \draw[gray,thin] (0.43,-0.76) -- (1.03,-0.34);
  \draw[gray,thin] (0.64,-1.10) -- (1.24,-0.68);
  \node[font=\scriptsize,gray] at (1.35,-1.10) {$\lambda_2$};

  \draw[->,thin] (0,0.62) arc (90:151:0.62);
  \node[font=\small] at (-0.45,0.72) {$i$};
  \draw[->,thin] (0,-0.55) arc (-90:-57:0.55);
  \node[font=\small] at (0.36,-0.65) {$r$};

  \node[font=\small,blue] at (1.72,-0.56) {$r<i$，靠近法线};
  \node[font=\small,align=center] at (2.0,0.92) {波前间距变小\\表示光速变慢};
  \draw[->,orange!80!black,thick] (1.40,0.55) -- (0.73,-0.08);
\end{tikzpicture}


\figcaption{折射率与光速：进入折射率更大的介质时，光速变小，光线靠近法线}\label{fig:refraction-speed-wavefront}
\end{center}

\textbf{频率、波长与颜色：}
光也是波，满足
\[
v=\lambda f.
\]
同一束光从一种介质进入另一种介质时，频率 $f$ 由光源决定，通常不变；传播速度 $v$ 改变，所以波长 $\lambda$ 随之改变。进入折射率更大的介质时，$v$ 变小，$\lambda$ 也变短。可见光的颜色主要由频率决定，所以红光、蓝光在同一介质中速度和波长可以不同，但“颜色身份”对应的是频率。

\textbf{认知提醒：} 判断波长变化时先问“是否换介质”。同一束光跨介质传播时，频率通常不变，先由介质判断速度变化，再由 $v=\lambda f$ 判断波长变化；只有在同一种介质中比较不同色光时，才把频率差异作为主要线索。
\end{conceptbox}

\begin{conceptbox}{全反射}

当光从光密介质射向光疏介质，且入射角大于临界角时，在理想无吸收界面的高中几何光学模型中不再画折射光，反射光全部返回原介质，这就是全反射。

\[
\boxed{\sin C=\frac{n_2}{n_1}}\qquad(n_1>n_2)
\]

\begin{center}
% TikZ 图：全反射和临界角
% Day 25：光学基础
\begin{tikzpicture}[scale=0.9,>=Latex]
  \fill[blue!7] (-3,-1.5) rectangle (3,0);
  \draw[thick] (-3,0) -- (3,0);
  \draw[dashed] (0,-1.5) -- (0,1.6) node[above] {法线};
  \node at (-2.25,-1.15) {光密介质 $n_1$};
  \node at (-2.25,0.55) {光疏介质 $n_2$};
  \draw[->,red,thick] (-1.6,-1.25) -- (0,0);
  \draw[->,red,thick] (0,0) -- (2.25,0);
  \draw[->,blue,thick] (-2.2,-1.05) -- (-0.1,0);
  \draw[->,blue,thick] (0.1,0) -- (2.2,-1.05);
  \draw (-90:0.65) arc (-90:-142:0.65);
  \node at (-116:0.95) {$C$};
  \node[red,above] at (1.45,0.10) {临界折射光};
  \node[red,above] at (1.65,0.45) {沿界面};
  \node[blue] at (1.55,-1.1) {全反射};
  \node[font=\small,align=left] at (4.95,-0.20)
    {$\sin C=\dfrac{n_2}{n_1}$\\[3pt]
     条件：$n_1>n_2$，且 $i>C$};
\end{tikzpicture}


\figcaption{全反射：只有从光密介质射向光疏介质才可能发生}\label{fig:total-internal-reflection}
\end{center}

光纤通信正是利用全反射，让光在细长光纤中多次反射而不轻易漏出。
\end{conceptbox}

\begin{conceptbox}{色散与光谱}

白光通过棱镜后会分成不同颜色的光，称为\textbf{色散}。原因是同一介质对不同频率的光折射率不同：通常玻璃对紫光的折射率较大、对红光的折射率较小，所以紫光偏折更明显。

\begin{center}
% TikZ 图：棱镜色散
% Day 25：光学基础
\begin{tikzpicture}[scale=0.9,>=Latex]
  \coordinate (A) at (0,0);
  \coordinate (B) at (2.0,0);
  \coordinate (C) at (1.0,1.75);
  \draw[thick,fill=blue!6] (A)--(B)--(C)--cycle;
  \draw[->,thick] (-2.2,0.9) -- (0.52,0.9) node[midway,above] {白光};
  \draw[thick] (0.52,0.9) -- (1.61,0.68);
  % 棱镜底边在下方：各色光均向底边偏折，紫光偏折最多。
  \draw[red,thick,->] (1.61,0.68) -- (3.7,0.12) node[right] {红};
  \draw[orange,thick,->] (1.61,0.68) -- (3.7,-0.08);
  \draw[green!60!black,thick,->] (1.61,0.68) -- (3.7,-0.28);
  \draw[blue,thick,->] (1.61,0.68) -- (3.7,-0.48);
  \draw[violet,thick,->] (1.61,0.68) -- (3.7,-0.70) node[right] {紫};
  \node[font=\small,align=center] at (1.0,-0.45) {不同颜色光折射率不同};
\end{tikzpicture}


\figcaption{棱镜色散：不同颜色光在同一介质中的折射率不同}\label{fig:prism-dispersion}
\end{center}

色散后的颜色按波长或频率排开，形成\textbf{光谱}。光谱并非只是一条彩带：不同物质的原子、分子具有不同的特征谱线，科学家可据此分析恒星大气和材料成分。
\end{conceptbox}

\begin{samepage}
\section*{四、公式与适用条件}

\textbf{作图前检查：}
\nopagebreak[4]
\begin{enumerate}
  \item 先判断是反射、折射、全反射还是干涉，不同模型不能混用。
  \item 光路可逆，但仍要标清入射光、反射光或折射光的传播方向。
  \item 色散题要比较不同颜色光在同一介质中的折射率；不要把颜色变化误认为频率在界面处改变。
  \item 干涉题先确认两束光相干，再判断路程差和条纹间距。
\end{enumerate}
\end{samepage}

\begin{center}
\begingroup
\renewcommand{\arraystretch}{1.18}
\begin{tabular}{p{0.24\textwidth}p{0.24\textwidth}p{0.42\textwidth}}
\toprule
公式 & 物理意义 & 适用条件 \\
\midrule
$i=i'$ & 反射定律 & 镜面反射；两角均相对于法线 \\
像距 $=$ 物距 & 平面镜成像 & 像与物关于镜面对称 \\
$n_1\sin i=n_2\sin r$ & 折射定律 & 两均匀介质分界面；角度均相对于法线 \\
$n=\dfrac{c}{v}$ & 绝对折射率 & $v$ 为该介质中的光速，$c$ 为真空光速 \\
$v=\lambda f$ & 光速、波长、频率关系 & 在同一介质内取同一束光的 $v,\lambda,f$；跨介质时 $f$ 通常不变 \\
$\sin C=\dfrac{n_2}{n_1}$ & 临界角 & 光密到光疏，$n_1>n_2$ \\
$\Delta \ell=k\lambda$ & 双缝亮纹条件 & 两缝为同相相干源，$k=0,\pm1,\pm2,\ldots$ \\
$\Delta \ell=(k+\frac12)\lambda$ & 双缝暗纹条件 & 两缝为同相相干源，$k=0,\pm1,\pm2,\ldots$ \\
$\Delta x=\dfrac{L\lambda}{d}$ & 双缝条纹间距 & 单色相干光；$L\gg d$ 且满足小角度近似 \\
\bottomrule
\end{tabular}
\endgroup
\end{center}

\begin{learningbox}{光学题的三类模型分析清单}
光学题先判断模型类型，再选择公式和作图方式：
\begin{enumerate}
  \item \textbf{界面光路：} 先画界面和法线，角度一律相对于法线量；反射用 $i=i'$，折射用 $n_1\sin i=n_2\sin r$，全反射还要先确认“光密到光疏”。
  \item \textbf{色散与光谱：} 先看不同颜色光在介质中的偏折程度；折射率较大的光更靠近法线，棱镜中通常紫光偏折更明显。
  \item \textbf{波动光学：} 先看是干涉、衍射还是偏振；双缝题先判断路程差的明暗条件，再用 $\Delta x=\dfrac{L\lambda}{d}$，计算前先统一纳米、毫米、米等单位。
\end{enumerate}
\textbf{简要归纳：} 几何光学先画光路，色散题先比较折射率，波动光学先找控制条纹或图样的量。

\textbf{来源理解：} 光路图负责整理角度和距离等几何关系，反射定律、折射定律和干涉条件则给出这些几何量必须满足的物理规律。图和定律缺一不可；若法线、界面、双缝、屏或光线方向没有标清，先不要急着代公式。
\end{learningbox}

\section*{五、典型模型}

\subsection*{模型 1：折射方向判断}

判断折射方向时，关键参照线是法线，而不是纸面上“向上弯、向下弯”的直观印象：
\[
n_2>n_1\Rightarrow r<i,\qquad n_2<n_1\Rightarrow r>i.
\]

\subsection*{模型 2：平面镜成像与视深}

平面镜成像和水中物体“看起来变浅”都不是新的力学过程，而是光路反向延长后形成的几何结果：
\begin{itemize}
  \item \textbf{平面镜：} 像与物关于镜面对称，像距等于物距，成正立等大的虚像。
  \item \textbf{视深：} 从水中射向空气的光线折射后远离法线，人眼沿空气中光线反向延长，会把水中物体看成位于较浅处的虚像。
\end{itemize}

\begin{center}
\begin{tikzpicture}[scale=0.86,>=Stealth]
  \tikzset{
    ray/.style={very thick,blue,->},
    ext/.style={blue!55,dashed,thick},
    guide/.style={gray,dashed},
    label/.style={font=\small}
  }

  \begin{scope}[xshift=-3.75cm]
    \node[font=\small\bfseries] at (0,1.85) {平行玻璃板};
    \fill[cyan!10] (-0.75,-1.40) rectangle (0.75,1.40);
    \draw[thick] (-0.75,-1.40) -- (-0.75,1.40);
    \draw[thick] (0.75,-1.40) -- (0.75,1.40);
    \node[font=\scriptsize] at (0,1.12) {$n_2$};
    \node[font=\scriptsize] at (-2.10,1.12) {$n_1$};
    \node[font=\scriptsize] at (2.45,-0.92) {$n_1$};

    \coordinate (P) at (-0.75,-0.20);
    \coordinate (Q) at (0.75,0.05);
    \draw[ray] (-3.0,-1.20) -- (P);
    \draw[ray] (P) -- (Q);
    \draw[ray] (Q) -- (3.0,1.05);
    \draw[blue!45,dashed] (P) -- (3.0,1.51);
    \draw[guide] (-1.35,-0.20) -- (0.05,-0.20);
    \draw[guide] (0.15,0.05) -- (1.35,0.05);
    \draw[<->,red!75!black] (2.20,1.11) -- (2.36,0.76);
    \node[font=\scriptsize,red!75!black] at (2.78,1.12) {侧向位移};
    \node[font=\scriptsize] at (-1.02,0.05) {$i$};
    \node[font=\scriptsize] at (-0.40,0.06) {$r$};
    \node[font=\scriptsize] at (1.18,-0.18) {$i'=i$};
    \node[font=\scriptsize,align=center] at (0,-1.72)
      {两次折射后方向恢复平行\\但传播路径发生侧移};
  \end{scope}

  \begin{scope}[xshift=3.85cm]
    \node[font=\small\bfseries] at (0,1.85) {近轴视深};
    \fill[blue!4] (-2.55,0) rectangle (2.55,1.40);
    \fill[cyan!9] (-2.55,-1.55) rectangle (2.55,0);
    \draw[thick] (-2.55,0) -- (2.55,0);
    \node[font=\scriptsize] at (-2.08,1.10) {$n_2$};
    \node[font=\scriptsize] at (-2.08,-1.28) {$n_1$};

    \coordinate (O) at (0,-1.28);
    \coordinate (Op) at (0,-0.78);
    \coordinate (P) at (0.66,0);
    \coordinate (Q) at (0.82,0);
    \coordinate (E1) at (1.55,1.05);
    \coordinate (E2) at (1.92,1.05);
    \fill[red] (O) circle (1.7pt) node[below right,label] {物点 $O$};
    \fill[red!55] (Op) circle (1.4pt);
    \node[font=\scriptsize,red!70!black,below right] at (0.10,-0.83) {虚像 $O'$};

    \draw[blue,very thick] (O) -- (P);
    \draw[ray] (P) -- (E1);
    \draw[blue,very thick] (O) -- (Q);
    \draw[ray] (Q) -- (E2);
    \draw[ext] (P) -- (Op);
    \draw[ext] (Q) -- (Op);
    \draw[thick] (1.74,1.05) ellipse (0.24 and 0.12);
    \node[font=\scriptsize] at (1.74,1.32) {眼};

    \draw[guide] (0,-1.48) -- (0,0.20);
    \draw[<->,gray] (-0.92,0) -- (-0.92,-0.78)
      node[midway,left,font=\scriptsize] {$h'$};
    \draw[<->,gray] (-1.65,0) -- (-1.65,-1.28)
      node[midway,left,font=\scriptsize] {$h$};
  \end{scope}
\end{tikzpicture}


\figcaption{平面镜成像与视深：虚像来自反射光线或折射光线的反向延长}\label{fig:plane-mirror-apparent-depth}
\end{center}

\textbf{学习边界：} 高中阶段遇到平面镜和视深题，重点是作图判断：先画实际光线，再画反向延长线。虚像是反向延长线的交点，不是真实光线会聚的位置。

\subsection*{模型 3：全反射与光纤}

全反射必须同时满足两个条件：
\[
\text{光密介质}\rightarrow\text{光疏介质},\qquad i>C.
\]
若只满足入射角大，但方向是从光疏到光密，不会发生全反射。

\subsection*{模型 4：色散与光谱}

色散题不能只记“紫光偏折大”。先明确比较的是\textbf{同一介质中的不同颜色光}：通常紫光频率较高、波长较短，玻璃对它的折射率较大，因此通过棱镜时更靠近法线、总偏折更明显。

光谱把不同波长的光按顺序展开。连续光谱常对应高温致密物体的发光；原子气体可呈现特征谱线。高中阶段不要求由谱线反推复杂原子结构，但要知道“不同物质的光谱不同”是光谱分析的依据。

\subsection*{模型 5：双缝干涉与路程差}

光具有波动性。双缝实验中，两条狭缝由同一束单色光照亮，在两缝处可视为同相相干源；发出的两列光波频率相同、相位差稳定。屏上某点的明暗由两路光的路程差 $\Delta\ell$ 决定：
\[
\text{亮纹：}\Delta\ell=k\lambda,\qquad
\text{暗纹：}\Delta\ell=\left(k+\frac12\right)\lambda.
\]
相邻亮纹间距为
\[
\Delta x=\frac{L\lambda}{d}.
\]

\begin{center}
% TikZ 图：双缝干涉
% Day 25：波动光学入门
\begin{tikzpicture}[scale=0.86,>=Latex]
  \draw[thick] (0,-1.8) -- (0,-0.35);
  \draw[thick] (0,0.35) -- (0,1.8);
  \draw[thick] (5,-2.1) -- (5,2.1);
  \fill[black] (0,0.55) circle (1.5pt) node[left] {$S_1$};
  \fill[black] (0,-0.55) circle (1.5pt) node[left] {$S_2$};
  \foreach \y in {-1.5,-0.75,0,0.75,1.5}{
    \draw[red!70,thick] (5,\y) -- (5.35,\y);
  }
  \draw[dashed] (0,0.55) -- (5,0.75);
  \draw[dashed] (0,-0.55) -- (5,0.75);
  \draw[<->] (0.28,-0.55) -- (0.28,0.55) node[midway,right] {$d$};
  \draw[<->] (0,-2.25) -- (5,-2.25) node[midway,below] {$L$};
  \node[right] at (5.35,0.75) {亮纹};
  \node[align=left,font=\small] at (7.2,0.25)
    {相邻亮纹间距：\\[2pt]
     $\Delta x=\dfrac{L\lambda}{d}$};
\end{tikzpicture}


\figcaption{双缝干涉：条纹间距随波长和屏距增大，随缝距增大而减小}\label{fig:double-slit-interference}
\end{center}

这个公式可以用“路程差”直观理解。相邻亮纹对应的路程差相差一个波长 $\lambda$。当屏到双缝距离 $L$ 远大于缝距 $d$ 时，小角度近似下
\[
\Delta \ell \approx d\frac{x}{L}.
\]
令相邻亮纹的路程差变化为 $\lambda$，可推出
\[
\Delta x=\frac{L\lambda}{d}.
\]
因此，屏越远或波长越长，条纹越疏；双缝间距越大，条纹越密。

\begin{center}
\begin{tikzpicture}[scale=0.95,>=Stealth]
  \draw[thick] (-3.2,-1.35) -- (-3.2,1.35);
  \draw[fill=white] (-3.2,0.45) circle (0.045);
  \draw[fill=white] (-3.2,-0.45) circle (0.045);
  \node[font=\small,left] at (-3.25,0.45) {$S_1$};
  \node[font=\small,left] at (-3.25,-0.45) {$S_2$};
  \draw[<->] (-3.55,-0.45) -- (-3.55,0.45);
  \node[font=\small,left] at (-3.55,0) {$d$};

  \draw[thick] (3.0,-1.55) -- (3.0,1.55);
  \coordinate (P) at (3.0,0.95);
  \coordinate (O) at (3.0,0);
  \fill[red] (P) circle (0.045);
  \node[font=\small,right] at (P) {$P$};
  \node[font=\small,right] at (O) {$O$};
  \draw[<->] (3.35,0) -- (3.35,0.95);
  \node[font=\small,right] at (3.35,0.48) {$x$};

  \draw[blue,thick] (-3.2,0.45) -- (P);
  \draw[blue,thick] (-3.2,-0.45) -- (P);
  \draw[dashed] (-3.2,0) -- (O);
  \draw[<->] (-3.2,-1.25) -- (3.0,-1.25);
  \node[font=\small] at (-0.1,-1.48) {$L$};
  \draw[->,thin] (-1.8,0) arc (0:10:1.4);
  \node[font=\small] at (-1.1,0.22) {$\theta$};

  \node[font=\small,align=center] at (0,1.45) {远屏近似：$\Delta \ell \approx d\sin\theta \approx d\dfrac{x}{L}$};
  \node[font=\small,align=center] at (0,-0.62) {亮纹：$\Delta \ell=k\lambda$\\相邻亮纹差一个 $\lambda$};
\end{tikzpicture}


\figcaption{双缝路程差：用几何关系把路程差转化为屏上距离}\label{fig:double-slit-path-difference}
\end{center}

\subsection*{模型 6：衍射、偏振与分辨率}

光通过小孔、窄缝或掠过障碍物边缘时会发生\textbf{衍射}。孔径或缝宽越小、波长越长，$\lambda/a$ 越大，衍射越明显；当孔径或缝宽 $a$ 与波长同量级时尤其显著。因此任何光学仪器的细节分辨能力都有物理极限。\textbf{偏振}则说明光是横波：偏振片能选择光的振动方向，液晶显示、偏光太阳镜和许多光学传感器都在利用这一性质。

\begin{center}
\begin{tikzpicture}[scale=0.92,>=Stealth]
  \begin{scope}
    \node[font=\small] at (0,1.75) {衍射极限};
    \draw[thick] (-2.2,-1.25) -- (-2.2,-0.22);
    \draw[thick] (-2.2,0.22) -- (-2.2,1.25);
    \foreach \y in {-0.8,0,0.8}{
      \draw[blue,->] (-3.35,\y) -- (-2.35,\y);
    }
    \foreach \r in {0.45,0.85,1.25,1.65}{
      \draw[blue!65] (-2.2,0) arc (180:35:\r);
      \draw[blue!65] (-2.2,0) arc (180:325:\r);
    }
    \draw[thick] (1.65,-1.25) -- (1.65,1.25);
    \draw[red,very thick] (1.65,0) -- ++(0.55,0);
    \draw[red!60,thick] (1.65,0.42) -- ++(0.28,0);
    \draw[red!60,thick] (1.65,-0.42) -- ++(0.28,0);
    \node[font=\scriptsize,align=center] at (0,-1.55) {孔径有限，像点会展开\\显微镜有分辨率上限};
  \end{scope}

  \begin{scope}[xshift=5.3cm]
    \node[font=\small] at (0,1.75) {偏振};
    \draw[thick] (-1.25,-1.10) rectangle (-0.75,1.10);
    \foreach \y in {-0.75,-0.45,...,0.75}{
      \draw[gray] (-1.20,\y) -- (-0.80,\y+0.25);
    }
    \node[font=\scriptsize] at (-1.0,-1.38) {偏振片};
    \draw[->,blue,thick] (-3.0,0) -- (-1.35,0);
    \foreach \ang in {0,30,60,90,120,150}{
      \draw[blue!65] (-2.55,0) -- ++(\ang:0.32);
      \draw[blue!65] (-2.55,0) -- ++(\ang+180:0.32);
    }
    \draw[->,red,thick] (-0.65,0) -- (1.25,0);
    \draw[red,very thick] (0.15,-0.45) -- (0.15,0.45);
    \draw[red,very thick] (0.75,-0.45) -- (0.75,0.45);
    \node[font=\scriptsize,align=center] at (0.25,-1.38) {只保留某一振动方向};
  \end{scope}
\end{tikzpicture}


\figcaption{光的衍射与偏振：衍射限制分辨率，偏振体现光的横波性质}\label{fig:optical-diffraction-polarization}
\end{center}

\textbf{学习边界：} 高中阶段重点是用现象说明物理意义：衍射说明光具有波动性，偏振进一步说明光是横波；不要求用复杂公式计算衍射图样。

\section*{六、例题精讲}

\begin{examplebox}{折射角计算}
\textbf{题意分析：} 光从空气射入某玻璃，入射角 $i=30^\circ$，玻璃折射率 $n=1.5$，空气折射率取 1。求折射角。

\textbf{物理模型：} 折射定律。

\begin{center}
% TikZ 图：折射角计算例题
% Day 25：光学基础
\begin{tikzpicture}[scale=0.86,>=Latex,line cap=round,line join=round]
  \draw[thick] (-2.55,0) -- (2.55,0);
  \draw[dashed] (0,-1.35) -- (0,1.35) node[above] {法线};
  \node[font=\small] at (-1.85,0.35) {空气 $n_1=1$};
  \node[font=\small] at (1.85,-0.35) {玻璃 $n_2=1.5$};
  \draw[->,blue,thick] (-1.05,1.25) -- (0,0);
  \draw[->,red,thick] (0,0) -- (0.70,-1.18);
  \draw (90:0.55) arc (90:130:0.55);
  \node[blue,font=\small] at (108:0.86) {$30^\circ$};
  \draw (-90:0.62) arc (-90:-60:0.62);
  \node[red,font=\small] at (-74:0.94) {$r$};
  \node[font=\small,align=center] at (0,-1.65) {进入光密介质，折射光靠近法线};
\end{tikzpicture}

\end{center}

\textbf{解题过程：}
\[
n_1\sin i=n_2\sin r
\]
\[
1\times\sin30^\circ=1.5\sin r
\]
\[
\sin r=\frac{1}{3},\qquad r\approx19.5^\circ.
\]

\textbf{方法总结：} 折射题先判断光从光疏到光密还是光密到光疏，再预判靠近或远离法线；计算结果要和预判一致。
\end{examplebox}

\begin{examplebox}{全反射临界角}
\textbf{题意分析：} 光从折射率为 $1.50$ 的玻璃射向空气，求临界角。

\textbf{物理模型：} 光密介质到光疏介质的全反射。

\begin{center}
% TikZ 图：全反射临界角例题
% Day 25：光学基础
\begin{tikzpicture}[scale=0.86,>=Latex,line cap=round,line join=round]
  \path[use as bounding box] (-2.95,-1.40) rectangle (3.15,1.35);
  \fill[blue!7] (-2.70,-1.25) rectangle (2.70,0);
  \draw[thick] (-2.70,0) -- (2.70,0);
  \draw[dashed] (0,-1.25) -- (0,1.20) node[above] {法线};
  \node[font=\small] at (-1.80,-0.92) {玻璃 $n_1=1.50$};
  \node[font=\small] at (-1.80,0.45) {空气 $n_2=1$};
  \draw[->,red,thick] (-1.05,-1.05) -- (0,0);
  \draw[->,red,thick] (0,0) -- (2.25,0);
  \draw (-90:0.55) arc (-90:-135:0.55);
  \node[red,font=\small] at (-112:0.86) {$C$};
  \node[red,font=\small] at (1.35,0.28) {临界时沿界面};
\end{tikzpicture}

\end{center}

\textbf{解题过程：}
\[
\sin C=\frac{n_2}{n_1}=\frac{1}{1.50}=\frac23.
\]
因此
\[
C\approx41.8^\circ.
\]

\textbf{方法总结：} 临界角题先确认光从光密介质射向光疏介质；公式中分子是光疏介质折射率，分母是光密介质折射率。
\end{examplebox}

\begin{examplebox}{双缝干涉条纹间距}
\textbf{题意分析：} 双缝间距 $d=0.50\ \mathrm{mm}$，屏到双缝距离 $L=1.0\ \mathrm{m}$，用波长 $\lambda=600\ \mathrm{nm}$ 的光照射，求相邻亮纹间距。

\textbf{物理模型：} 双缝干涉。

\begin{center}
% TikZ 图：双缝干涉例题
% Day 25：光学基础
\begin{tikzpicture}[scale=0.80,>=Latex,line cap=round,line join=round]
  \draw[thick] (0,-1.25) -- (0,-0.54);
  \draw[thick] (0,-0.36) -- (0,0.36);
  \draw[thick] (0,0.54) -- (0,1.25);
  \fill[black] (0,0.45) circle (1.4pt) node[left] {$S_1$};
  \fill[black] (0,-0.45) circle (1.4pt) node[left] {$S_2$};
  \draw[thick] (5.00,-1.35) -- (5.00,1.35);
  \foreach \y in {-0.72,0,0.72} {
    \draw[red!70,thick] (5,\y) -- (5.35,\y);
  }
  \draw[<->] (0.25,-0.45) -- (0.25,0.45) node[midway,right] {$d=0.50\ \mathrm{mm}$};
  \draw[<->] (0,-1.45) -- (5,-1.45) node[midway,below] {$L=1.0\ \mathrm{m}$};
  \draw[<->,orange,thick] (5.55,0) -- (5.55,0.72) node[midway,right] {$\Delta x$};
  \node[font=\small,align=center] at (3.05,1.25) {$\lambda=600\ \mathrm{nm}$\\先统一单位};
\end{tikzpicture}

\end{center}

\textbf{解题过程：}
\[
\Delta x=\frac{L\lambda}{d}
=\frac{1.0\times600\times10^{-9}}{0.50\times10^{-3}}
=1.2\times10^{-3}\ \mathrm{m}.
\]
即
\[
\Delta x=1.2\ \mathrm{mm}.
\]

\textbf{方法总结：} 干涉题先统一到国际单位制，尤其纳米、毫米，再代入 $\Delta x=\dfrac{L\lambda}{d}$。
\end{examplebox}

\section*{七、变式训练}

\begin{enumerate}
  \item 光从水射入空气时，折射光线相对法线是靠近还是远离？
  \item 某介质中光速为 $2.0\times10^8\ \mathrm{m/s}$，求折射率。
  \item 光从折射率 $1.33$ 的水射向空气，写出临界角满足的方程。
  \begin{center}
  \begin{tikzpicture}[scale=0.82,>=Stealth]
  \draw[thick] (-2.4,0) -- (2.4,0);
  \draw[dashed] (0,-1.55) -- (0,1.6) node[above] {法线};
  \draw[->,blue,thick] (-1.7,-1.1) -- (0,0);
  \draw[->,blue,thick] (0,0) -- (1.8,0);
  \draw (0,-0.75) arc[start angle=-90,end angle=-147,radius=0.75];
  \node at (-0.48,-0.82) {$C$};
  \node at (-1.55,-0.35) {水};
  \node at (1.55,0.38) {空气};
  \node[font=\small] at (1.15,-1.18) {$\sin C=\dfrac{n_{\text{空气}}}{n_{\text{水}}}$};
\end{tikzpicture}

  \end{center}
  \item 白光通过玻璃棱镜后，红光和紫光哪一种偏折更明显？从折射率角度说明理由。
  \item 双缝实验中，某点两路光的路程差为 $3\lambda$，该点是亮纹还是暗纹？
  \item 双缝实验中，若缝距增大为原来的 2 倍，条纹间距怎样变化？
  \begin{center}
  \begin{tikzpicture}[scale=0.8,>=Stealth]
  \draw[thick] (0,-0.65) -- (0,-0.32);
  \draw[thick] (0,-0.18) -- (0,0.18);
  \draw[thick] (0,0.32) -- (0,0.65);
  \fill (0,0.25) circle (1.2pt);
  \fill (0,-0.25) circle (1.2pt);
  \node[left=6pt] at (0,0.25) {$S_1$};
  \node[left=6pt] at (0,-0.25) {$S_2$};
  \draw[<->] (0.25,-0.25) -- (0.25,0.25) node[midway,right] {$d$};
  \draw[thick] (2.4,-0.9) -- (2.4,0.9);
  \foreach \y in {-0.5,0,0.5} \draw[red] (2.4,\y) -- (2.6,\y);
  \draw[<->,orange,thick] (2.75,0) -- (2.75,0.5) node[midway,right] {$\Delta x$};
  \node at (1.2,-1.2) {原来};
  \begin{scope}[xshift=4.2cm]
    \draw[thick] (0,-0.85) -- (0,-0.57);
    \draw[thick] (0,-0.43) -- (0,0.43);
    \draw[thick] (0,0.57) -- (0,0.85);
    \fill (0,0.5) circle (1.2pt);
    \fill (0,-0.5) circle (1.2pt);
    \node[left=6pt] at (0,0.5) {$S_1$};
    \node[left=6pt] at (0,-0.5) {$S_2$};
    \draw[<->] (0.25,-0.5) -- (0.25,0.5) node[midway,right] {$2d$};
    \draw[thick] (2.4,-0.9) -- (2.4,0.9);
    \foreach \y in {-0.25,0,0.25} \draw[red] (2.4,\y) -- (2.6,\y);
    \draw[<->,orange,thick] (2.75,0) -- (2.75,0.25) node[midway,right] {$\Delta x/2$};
    \node at (1.2,-1.2) {缝距加倍};
  \end{scope}
\end{tikzpicture}

  \end{center}
\end{enumerate}

\begin{learningbox}{变式训练参考结果}
\begin{enumerate}[itemsep=0.45em,topsep=0.3em,parsep=0pt]
  \item 光从水进入空气，即从光密介质进入光疏介质，折射光线远离法线。
  \item $n=c/v=1.50$。
  \item 临界角满足 $\sin C=n_{\text{空气}}/n_{\text{水}}=1/1.33$，
  即 $C\approx48.8^\circ$。
  \item 紫光偏折更明显。同一玻璃对紫光的折射率通常大于对红光的折射率，紫光在玻璃中的传播速度更小。
  \item 路程差 $3\lambda$ 是波长的整数倍，两列相干光同相叠加，该点为亮纹。
  \item 由 $\Delta x=L\lambda/d$，缝距增大为原来的 2 倍时，条纹间距变为原来的一半。
\end{enumerate}
\end{learningbox}

\section*{八、竞赛入门补充}
\begin{competitionbox}
\textbf{费马原理的思想}

几何光学中有一个更深的统一观点：光在两点之间传播时，实际光路使传播时间取驻值，也就是相对于邻近路径的一阶变化为零。反射定律和折射定律都可以从这个思想推出。

\textbf{为什么折射会偏折？}
光在不同介质中速度不同。若一束斜入射光进入较慢介质，波前的一侧先进入介质而速度先变慢，另一侧仍在原介质中较快传播，波前发生旋转，于是光线方向改变。

\textbf{竞赛入门模型：} 多介质光路问题通常先把光路分成几段，分别写出每段时间
\[
t=\frac{s_1}{v_1}+\frac{s_2}{v_2}+\cdots,
\]
再利用几何关系寻找传播时间取驻值的路径。高中阶段先把法线、角度和折射率用清楚，竞赛阶段再把它们统一到费马原理的思想中。
\end{competitionbox}

\subsection*{1.\quad 波动光学的三个关键词}

\begin{enumerate}
  \item \textbf{相干：} 两束光频率相同、相位差稳定，才容易形成稳定干涉条纹。
  \item \textbf{路程差：} 在空气中的双缝实验里，决定某处是加强还是减弱。
  \item \textbf{波长：} 条纹间距、衍射明显程度都与波长有关。
\end{enumerate}

双缝干涉中，相邻亮纹间距
\[
\Delta x=\frac{L\lambda}{d}
\]
说明：波长越长、屏越远，条纹越疏；双缝间距越大，条纹越密。

\textbf{公式来源直觉：} 屏越远，同样的角度差在屏上拉开的距离越大；双缝间距越大，同一位置产生的路程差变化更快，条纹就更密。这比只记公式更能防止把 $L$ 和 $d$ 的影响记反。

\subsection*{2.\quad 薄膜干涉：从光程差到彩色条纹}

肥皂泡、油膜和镜头镀膜常出现彩色条纹，原因是\textbf{薄膜干涉}。入射光到达薄膜上表面时，一部分直接反射；另一部分进入薄膜，在下表面反射后再射出。两束反射光沿同一观察方向传播时可以叠加。

\begin{center}
\begin{tikzpicture}[scale=0.92,>=Stealth]
  \tikzset{
    ray/.style={very thick,blue,->},
    rayb/.style={very thick,orange!90!black,->},
    guide/.style={gray,dashed},
    film/.style={fill=cyan!10,draw=cyan!60!black,thick},
    note/.style={font=\small,align=center}
  }

  \fill[film] (-3.35,-0.55) rectangle (3.35,0.55);
  \draw[thick] (-3.35,0.55) -- (3.35,0.55);
  \draw[thick] (-3.35,-0.55) -- (3.35,-0.55);
  \node[font=\small] at (-2.80,0.92) {空气};
  \node[font=\small,cyan!60!black] at (-2.75,0) {薄膜};
  \node[font=\small] at (-2.80,-0.92) {空气/玻璃};

  \coordinate (P) at (-0.70,0.55);
  \coordinate (Q) at (0.05,-0.55);
  \coordinate (R) at (0.82,0.55);

  \draw[ray] (-2.55,1.45) -- (P);
  \draw[ray] (P) -- (-1.68,1.42);
  \node[font=\scriptsize,blue] at (-1.92,1.62) {反射光 1};

  \draw[rayb] (P) -- (Q);
  \draw[rayb] (Q) -- (R);
  \draw[rayb] (R) -- (1.84,1.42);
  \node[font=\scriptsize,orange!90!black] at (2.05,1.62) {反射光 2};

  \draw[guide] (P) -- ++(0,-1.32);
  \draw[guide] (Q) -- ++(0,1.32);
  \draw[<->,gray] (3.00,-0.55) -- (3.00,0.55) node[midway,right,font=\small] {$d$};
  \node[note] at (0,-1.48)
    {两束反射光来自薄膜上下表面，光程差约与 $2nd$ 有关};
  \node[note] at (0,-1.95)
    {不同颜色满足加强或减弱条件不同，于是出现彩色条纹};
\end{tikzpicture}


\figcaption{薄膜干涉：上下表面的两束反射光叠加，光程差决定加强或减弱}\label{fig:thin-film-interference}
\end{center}

\textbf{核心判断：} 薄膜厚度 $d$ 改变时，两束反射光的光程差随之改变。不同波长的光满足加强或减弱条件的位置不同，所以白光照射下会看到彩色条纹。

\textbf{模型边界：} 薄膜干涉的严格明暗条件还要考虑反射时是否发生半波损失。高中入门阶段先抓住“同一束光分成两束，再因光程差发生干涉”的模型；遇到定量题时再按题目给出的相位条件处理。

\textbf{科学应用：} 镜头增透膜会设计合适厚度，使某些反射光相互削弱，从而减少反光、提高透光率。这里用到的仍然是干涉思想，而不是简单“涂一层透明材料”。

\begin{examplebox}{光纤临界角与数值判断}
\textbf{题意分析：} 某光纤芯层折射率 $n_1=1.48$，包层折射率 $n_2=1.46$。光从芯层射向包层，求临界角满足的关系，并判断临界角接近 $80^\circ$ 还是 $40^\circ$。

\textbf{物理模型：} 光密介质到光疏介质，全反射临界角。

\textbf{解题过程：}
\[
\sin C=\frac{n_2}{n_1}=\frac{1.46}{1.48}\approx0.986.
\]
由于 $\sin C$ 接近 1，所以 $C$ 接近 $90^\circ$；并且 $\sin80^\circ\approx0.985$，与 $0.986$ 很接近，而 $\sin40^\circ\approx0.643$ 明显偏小，因此临界角更接近 $80^\circ$。

\textbf{方法总结：} 全反射只能从光密介质射向光疏介质时发生。折射率差很小时，临界角很大，入射角必须更接近 $90^\circ$、光线更“贴近界面”，才可能满足 $i>C$。
\end{examplebox}

\subsection*{3.\quad 分层训练}

\begin{enumerate}
  \item \textbf{基础：} 光从空气进入水中，入射角为 $45^\circ$，折射角比入射角大还是小？
  \item \textbf{提高：} 某透明介质折射率为 $1.60$，光从该介质射向空气，写出临界角方程。
  \item \textbf{综合：} 双缝实验中换用波长更短的蓝光，其他条件不变，条纹间距如何变化？说明理由。
  \item \textbf{拓展：} 肥皂泡为什么会呈现彩色条纹？从薄膜干涉角度说明。
  \item \textbf{开放思考：} 为什么光学仪器的分辨能力会受到衍射限制？
\end{enumerate}

\section*{九、易错点}
\begin{mistakebox}
\begin{enumerate}
  \item \textbf{角度量错：} 作图第一步先在入射点画法线，再从法线量入射角、反射角和折射角；若题目给的是“与界面夹角”，要先用 $90^\circ$ 减去它。
  \item \textbf{把虚像当成真实光线交点：} 先区分实线光路和虚线延长线。平面镜像和视深中的像都来自光线反向延长，真实光线并没有在虚像处会聚。
  \item \textbf{全反射条件不完整：} 先判断是否光密到光疏，再看入射角是否满足 $i>C$。
  \item \textbf{把折射率理解成光速变快：} 先用 $n=c/v$ 判断；折射率越大，介质中光速越小。
  \item \textbf{误以为光进入介质后频率改变：} 跨介质题先固定频率。对同一束光，频率通常不变，改变的是速度和波长。
  \item \textbf{把色散当作频率改变：} 光在同一界面折射时频率通常不变；不同颜色光的偏折差异来自介质对它们的折射率不同。
  \item \textbf{干涉题不检查相干条件：} 稳定条纹要求两束光频率相同、相位差稳定；双缝实验用同一束光照亮两缝，正是为了满足这一条件。
  \item \textbf{波动光学单位混乱：} 计算前先把 $\lambda$、$d$、$L$ 统一到米；$\lambda$ 常用纳米，$d$ 常用毫米，是最容易漏换的两个量。
\end{enumerate}
\end{mistakebox}

\section*{十、今日小结}
\begin{enumerate}
  \item 反射角等于入射角，折射遵循 $n_1\sin i=n_2\sin r$。
  \item 平面镜成像满足像距等于物距；视深变浅来自折射光线的反向延长。
  \item 折射率 $n=c/v$，折射率越大，介质中光速越小。
  \item 光满足 $v=\lambda f$；跨介质传播时频率通常不变，速度和波长改变。
  \item 全反射必须从光密介质射向光疏介质，且入射角大于临界角。
  \item 白光经棱镜发生色散；通常紫光折射率较大、偏折更明显，光谱可用于分析物质成分。
  \item 双缝干涉说明光具有波动性，条纹间距 $\Delta x=L\lambda/d$。
  \item 薄膜干涉来自上下表面反射光的叠加，可解释肥皂泡彩纹和增透膜。
  \item 衍射限制光学仪器分辨率，偏振体现光的横波性质。
\end{enumerate}

\textbf{一分钟复盘：} 任选一道光学题，先画法线或双缝与屏，再标角度、介质、光线方向或路程差；确认模型后再写公式。
